\documentclass[11pt]{article}
% Internal note (not part of any submitted paper).
\usepackage[letterpaper,margin=25mm]{geometry}
\usepackage{amsmath,amssymb,amsthm,bm}
\usepackage{booktabs}
\usepackage{algorithm}
\usepackage[noend]{algpseudocode}
\usepackage[round,authoryear]{natbib}
\usepackage{url}
\usepackage[hidelinks]{hyperref}

\newtheorem{lemma}{Lemma}
\newtheorem{theorem}{Theorem}
\newcommand{\Xset}{\mathcal X}
\newcommand{\E}{\mathbb E}

\title{Scoring column pairs under a $K$-class profile mixture\\
for statistical progressive alignment}
\author{Internal note, tkf-dp}
\date{September 2026}

\begin{document}
\maketitle

\begin{abstract}
In progressive alignment under a single substitution model, scoring the merge
of a column of one clade's alignment with a column of its sibling's is a dot
product of two length-$A$ vectors of partial likelihoods; under a $K$-class
profile mixture such as C20 it is a dot product of length $KA$. We derive a
variational lower bound whose per-cell cost is $O(K+A)$, built from fixed-size
column summaries by a soft version of Fitch's rule, and find that at the
divergences of BAliBASE it is too loose to drive alignment decisions. A
sparse-class lower bound on the exact score is practically exact at half its
cost, and collapsing the classes of an exact posterior identity gives an $O(A)$
score that is as accurate for alignment. With one class this score is exact
and has the form of MUSCLE's log-expectation score. We compare the scores as
bounds, as alignment scores and as MCMC proposals; characterize Viterbi,
hybrid and posterior decoding under MixFrag, whose Viterbi path is
systematically unfair to fragment mixtures; and benchmark a JAX implementation
with expected-SPS decoding on BAliBASE, where it reaches SP 0.751 against 0.841
for MUSCLE.
\end{abstract}

\section{Introduction}
\label{sec:intro}

\paragraph{Progressive multiple alignment.} CLUSTAL \citep{HigginsSharp1988}
established the progressive strategy that most multiple aligners still
follow: build a guide tree from pairwise similarities, then align clusters of
sequences pairwise, following the branching order of the tree, so that each
step aligns two profiles. CLUSTAL~W \citep{Thompson1994} added sequence
weighting, position-specific gap penalties and stage-dependent substitution
matrices. MUSCLE \citep{Edgar2004} replaced the guide-tree distances by fast
$k$-mer distances and scored two profile columns with a log-expectation
function built from the residue frequencies of each column, so that the
substitution score adapts to the composition of each position. ProbCons
\citep{Do2005} scored alignments with pair-HMM posterior probabilities made
consistent across sequences, and reported higher accuracy than CLUSTAL~W,
T-Coffee and MUSCLE on all but one of the BAliBASE reference sets; its extended variant, ProbCons-ext, added a
second pair of insertion states for long or terminal insertions, i.e.\ a
two-component mixture of gap lengths. Position-specific substitution scoring
and mixtures of gap lengths are the two ingredients this note brings together
within a single evolutionary model.

\paragraph{Statistical alignment.} The interpretation of alignment scores in
terms of a likelihood model of sequence evolution has been called
\emph{statistical alignment} \citep{Hein2000}. Its pair HMMs come from
continuous-time models of insertion and deletion: TKF91
\citep{ThorneKishinoFelsenstein1991}, in which indels are single residues;
TKF92 \citep{ThorneKishinoFelsenstein1992}, in which they are fragments of
geometrically distributed length; and MixFrag \citep{HolmesMixFrag}, in which
each fragment draws its extension probability from a finite mixture, so that
gap lengths follow a mixture of geometric distributions. The substitution
process is a continuous-time Markov chain (CTMC) on the $A$ residues, such as
LG \citep{LeGascuel2008}. In a progressive statistical aligner each column of
a clade's MSA is summarized by its vector of Felsenstein partial likelihoods
\citep{Felsenstein81}, and scoring a pair of columns from sibling clades is a
dot product of two such vectors, of length $A$.

\paragraph{Position-specific scoring by mixtures.} The natural likelihood
analogue of position-specific scoring is a mixture over columns: replace the
single CTMC by $K$ CTMCs that share exchangeabilities but differ in their
equilibrium profile (and possibly rate), with each column drawing its class
independently. Empirical profile mixtures such as C20 and C60
\citep{LeGascuelLartillot2008}, usually combined with LG exchangeabilities,
have exactly this form. CAT \citep{Lartillot2004} is the nonparametric
version, a Dirichlet-process mixture over profiles in which the number of
classes is inferred; it fits the same form once its profiles are fixed, for
example at a posterior sample. The cost is in the partial likelihoods: each
column now carries $K$ vectors of length $A$, scoring two columns is a dot
product of length $KA$, and each profile--profile alignment of MSAs with
$L_A$ and $L_B$ columns costs $O(KAL_AL_B)$, for example $20$ to $60$ times
the single-CTMC cost.

\paragraph{This note.} We develop a variational lower bound on each column's
likelihood whose profile--profile step costs $O((K+A)L_AL_B)$
(Section~\ref{sec:bound}). It takes an expectation of the complete-data
log-likelihood under a single distribution over substitution histories that is
shared by all classes, so each class's contribution is a fixed linear function
of a few expected statistics, and the per-class terms of a merge separate into
a length-$K$ dot product, a length-$A$ dot product and a sparse correction
(Section~\ref{sec:separates}). The history distribution is built bottom-up by a
rule that is \emph{Fitch's rule made soft}. At the divergences of BAliBASE the
bound turns out to be too loose for alignment decisions
(Section~\ref{sec:tightness}); a sparse-class bound and a class-collapsed
posterior score repair this, and the latter has the form of MUSCLE's
log-expectation score. We compare all the scores for likelihoods, decoding and
MCMC proposals (Section~\ref{sec:choose}), analyse the decoding objectives
(Section~\ref{sec:decoding}), and report a benchmark
(Section~\ref{sec:results}).

\paragraph{Relation to Fitch and Sankoff parsimony.} Fitch's algorithm
\citep{Fitch1971} computes, in one pass from the leaves up, a set of candidate
states at each internal node, the intersection of the children's sets if it is
nonempty and otherwise their union, and counts one change whenever it takes
the union; Sankoff's algorithm \citep{Sankoff1975} generalizes this to
arbitrary substitution costs by dynamic programming over states. In our
construction (Section~\ref{sec:summaries}), each internal node instead carries
a \emph{distribution} over states. When two children merge, their states
are drawn independently from their distributions, and the path joining them
through the parent has no substitution if the two states agree and exactly one,
at a uniformly random point, if they differ. So the parent inherits the shared
state when the children agree (a soft intersection); when they disagree it
inherits one child's state or the other's (a soft union), favouring the child
closer to it in proportion to the other child's share of the path length. The
expected number of changes grows by the probability that the two children
disagree, a soft Fitch count. There are three differences from parsimony.
\begin{itemize}
\item \emph{Costs come from the model.} Each change $a\to b$ on a branch of
  length $t$ contributes $\log S_{ab}+\log t$, and time spent in each state is
  charged by the exit rates, so the implied costs are those of weighted
  (Sankoff) parsimony with model-derived, branch-length-dependent weights,
  evaluated in expectation rather than minimized.
\item \emph{There is no downward pass.} Fitch and Sankoff resolve ancestral
  states with a second pass from the root; here each node sees only the leaves
  below it, and the two children of a merge are independent under $q$. This
  independence is the main source of looseness in the bound
  (Section~\ref{sec:pfam}); re-estimating each clade's summaries from the
  exact posterior (Section~\ref{sec:reest}) is the probabilistic counterpart of
  the downward pass, and is done once per column outside the dynamic program.
\item \emph{It defines a path measure.} $q$ is a probability measure on
  continuous-time substitution histories, not a labelling of nodes: a change
  on a branch occurs at a time drawn uniformly along it. Densities are taken
  with respect to counting measure on discrete configurations times Lebesgue
  measure on the substitution times (Section~\ref{sec:bound}), so the bound
  includes the entropy of these times. A single most-parsimonious history is
  the degenerate case in which $q$ concentrates on one discrete configuration,
  giving the weaker bound $\log p(D)\ge\log p(h,D)$. As branch lengths shrink,
  each extra change costs a factor of $t$, so the likelihood is dominated by
  minimum-change histories: the regime in which parsimony and likelihood agree,
  and which a Fitch-like construction targets.
\end{itemize}

\paragraph{Outline.} Section~\ref{sec:model} states the model,
Section~\ref{sec:bound} the bound, Section~\ref{sec:summaries} the summaries
and the merge rule, and Section~\ref{sec:separates} the separable pair score.
Section~\ref{sec:reest} replaces the soft-Fitch summaries inside each clade by
exact posterior statistics, Section~\ref{sec:pfam} tests the bound on Pfam,
Section~\ref{sec:viterbi} gives pseudocode for Viterbi progressive
alignment under MixFrag, with the expected-statistics computation in
Appendix~\ref{app:hr}, and Section~\ref{sec:distances} the pairwise distances
for the guide tree, with an alignment-free initialization in
Appendix~\ref{app:kgram}. Section~\ref{sec:tightness} shows that the bound is too
loose for alignment decisions at BAliBASE divergences and derives the
alternatives, Section~\ref{sec:choose} compares all the scores,
Section~\ref{sec:decoding} the decoding objectives, and
Section~\ref{sec:results} reports the benchmark.

\section{Model}
\label{sec:model}

A column has leaf states $D$ on a tree with branch lengths, over an alphabet
$\Xset$ of size $A$. Columns evolve independently. Column $c$ has a latent
class $k\in\{1,\dots,K\}$ with prior weight $w_k$; given $k$, residues evolve
under the reversible generator
\[
Q^k_{xy}=r_k\,S_{xy}\,\pi^k_y\quad(x\ne y),\qquad
q^k_x=-Q^k_{xx}=r_k\,(S\pi^k)_x ,
\]
with exchangeabilities $S=S^\top\ge0$ ($S_{xx}=0$) shared by all classes, and a
class-specific profile $\pi^k$ and rate $r_k$. Profile mixtures such as CAT
\citep{Lartillot2004} and C20/C60 \citep{LeGascuelLartillot2008} have this form. The column
likelihood is $p(D)=\sum_kw_k\,p_k(D)$, where $p_k(D)=p(D\mid\pi^k,r_k)$.

\paragraph{Exact pair score.} Let $F^k_A(x)$ and $F^k_B(x)$ be the partial
likelihoods of column $m$ of clade $A$ and column $n$ of clade $B$ at the
clades' top nodes, and $M^k(t)=e^{tQ^k}$. By reversibility the two branches
to the parent combine into one of length $t=t_A+t_B$, and the exact score of
the merged column is
\[
\log\sum_kw_k\sum_x\big(\pi^k_xF^k_A(x)\big)\big(M^k(t)F^k_B\big)(x).
\]
Only one side needs pre-multiplying by a transition matrix, once per column
and outside the DP, at $O(KA^2)$ per column. Each DP cell is then a dot
product of length $KA$, and each column stores $KA$ partials.

\section{A bound linear in the column summaries}
\label{sec:bound}

\paragraph{Paths.} A \emph{path} $h$ is a realization of the chain on every
branch: the root state and, on each branch, the times and targets of all
substitutions. Its statistics are $V'_x(h)$, the number of \emph{tokens} equal
to $x$ (the root state plus the destination of every substitution);
$U(h)=\sum_xV'_x(h)-1$, the number of substitutions; $T_x(h)$, the total time
spent in state $x$ summed over branches; and
$g(h)=\prod_{\text{substitutions }x\to y}S_{xy}$.

\paragraph{Path space.} A path consists of a discrete \emph{configuration}
(the root state and, for each branch $e$, the number $n_e$ of substitutions and
their target states) and the substitution times, which lie in
$\prod_e\{0<\tau_1<\dots<\tau_{n_e}<t_e\}\subset\mathbb R^{U(h)}$ with
$U(h)=\sum_en_e$. Path space is therefore a disjoint union of pieces of
different dimensions, one per configuration. The reference measure is counting
measure on configurations times $U(h)$-dimensional Lebesgue measure on the
times; a configuration with no substitutions is a point of mass one. All
densities and entropies below are taken with respect to this measure, so their
dependence on the time unit cancels in the bound.

\begin{lemma}[path density]\label{lem:density}
With respect to this reference measure,
\[
\log p_k(h)=\sum_xV'_x(h)\log\pi^k_x+U(h)\log r_k+\log g(h)
 -r_k\sum_xT_x(h)(S\pi^k)_x .
\]
\end{lemma}
\begin{proof}
The standard density of a continuous-time Markov chain path is the root
probability times the product of the rates of the jumps taken times
$\exp(-\int q^k_{X(u)}\,du)$. The root contributes $\log\pi^k_{x_0}$, each jump
$x\to y$ contributes $\log r_k+\log S_{xy}+\log\pi^k_y$, and the exponent is
$\sum_xT_xq^k_x$.
\end{proof}

\begin{lemma}[re-rooting]\label{lem:reroot}
Because each class is reversible, $p_k(D)$ does not depend on where the root
is placed \citep{Felsenstein81}. We may root the tree at any node, including a
leaf.
\end{lemma}

\begin{theorem}[per-class bound]\label{thm:bound}
Let $q$ be any distribution over paths that agrees with $D$ at the leaves, with
a density (with respect to the same reference measure) whose support lies in
that of $p_k$. Write $V'=\E_q[V'(h)]$, $U=\E_q[U(h)]$ and $T=\E_q[T(h)]$ for
its expected statistics, and
$C=\E_q[\log g(h)]+H(q)$. Then
\begin{equation}\label{eq:Phi}
\log p_k(D)\ \ge\ \Phi_k(s)=C+\sum_xV'_x\log\pi^k_x+U\log r_k
 -r_k\sum_xT_x(S\pi^k)_x ,
\end{equation}
where $s=(V',U,T,C)$ is the \emph{summary} of $q$. Consequently
\begin{equation}\label{eq:mix}
\log p(D)\ \ge\ \log\sum_kw_k\,e^{\Phi_k(s)} .
\end{equation}
\end{theorem}
\begin{proof}
Jensen's inequality gives
$\log p_k(D)=\log\E_q[p_k(h)/q(h)]\ge\E_q[\log p_k(h)]+H(q)$; substitute
Lemma~\ref{lem:density} and take expectations, which are linear in $V',U,T$.
Then \eqref{eq:mix} follows because $\log$ is increasing and
$p_k(D)\ge e^{\Phi_k(s)}$ for every $k$.
\end{proof}

The key property is that $\Phi_k$ is \emph{affine in the summary $s$}, and
that one $q$, and hence one summary, serves every class.

\section{Summaries that compose across merges}
\label{sec:summaries}

Each column of an MSA carries a summary $s=(V',U,T,C)$ together with
$q_{\mathrm{top}}$, the distribution under $q$ of the state at the clade's top
node: $O(A)$ numbers, independent of the depth of the MSA. A leaf with residue
$x$ has $V'=e_x$, $U=0$, $T=0$, $C=0$ and $q_{\mathrm{top}}=e_x$.

\begin{lemma}[merge]\label{lem:merge}
Let $q_A$ and $q_B$ be path distributions on the subtrees below sibling nodes
$A$ and $B$, and $q_{\mathrm{br}}(\cdot\mid a,b)$ a distribution of paths on
the branch of length $t=t_A+t_B$ joining $A$ to $B$ through their parent $P$.
Then $q=q_A\,q_B\,q_{\mathrm{br}}(\cdot\mid x_A,x_B)$ is a valid path
distribution for the merged tree, rooted where $A$'s tree is rooted, and
\begin{itemize}
\item $V'=V'_A+V'_B-\E[e_{x_B}]+\E[\text{destinations on the bridge}]$;
\item $U$, $T$ and $C$ add, with the bridge's contributions added in;
\item $q_{\mathrm{top}}(P)$ is the marginal of $q_{\mathrm{br}}$ at distance
  $t_A$ from $A$.
\end{itemize}
\end{lemma}
\begin{proof}
The product is a normalized distribution over paths that agree with $D$. By
the chain rule,
$H(q)=H(q_A)+H(q_B)+\E\,H\big(q_{\mathrm{br}}(\cdot\mid x_A,x_B)\big)$, and
$\log g$ adds over branches, so $C$ adds.
$B$'s summary was computed with $B$'s tree rooted at one of its leaves,
$\ell_B$; in the merged tree, $B$'s subtree hangs from node $B$. Only the
branches on the path between $\ell_B$ and $B$ change direction. Along that
path the visited states $s_0,\dots,s_j$ give the same multiset of tokens in
either direction (the start state plus the destinations, or the end state plus
the sources), so $B$'s own tokens, directed away from $B$, are
$V'_B-e_{x_B}$; node $B$'s state now arrives through the bridge. Dwell times,
the number of substitutions and $\log g$ (as $S$ is symmetric) do not depend
on direction, and time reversal is a measure-preserving bijection, so $U$,
$T$ and the entropy are unchanged. The next merge conditions on $x_P$, a
function of the bridge, so the construction repeats and needs only
$q_{\mathrm{top}}(P)$.
\end{proof}

\paragraph{Single-jump bridge.} Take $x_A\sim q_{A,\mathrm{top}}$ and
$x_B\sim q_{B,\mathrm{top}}$ independently; write $q_A,q_B$ for these two
distributions. If $x_A\ne x_B$, the bridge makes one substitution
$x_A\to x_B$ at a time uniform on $(0,t)$: it contributes $\log S_{x_Ax_B}$ to
$\log g$ and $\log t$ to the entropy (this needs $S_{x_Ax_B}>0$; otherwise the
bound is $-\infty$, which is still valid). If $x_A=x_B$, the bridge has no
substitutions and zero entropy. Taking expectations, the merge adds
\begin{align*}
\Delta V'&=q_B\circ(1-q_A)-q_B=-\,q_A\circ q_B, &
\Delta U&=1-q_A^\top q_B, \\
\Delta T&=\tfrac t2\,(q_A+q_B), &
\Delta C&=q_A^\top L\,q_B,\quad L=(\log S+\log t)\circ(1-I),
\end{align*}
where $\circ$ is the elementwise product, and
\[
q_{\mathrm{top}}(P)=q_A\circ q_B+\tfrac{t_B}{t}\,q_A\circ(1-q_B)
 +\tfrac{t_A}{t}\,q_B\circ(1-q_A).
\]
By Lemma~\ref{lem:merge} these are exact expectations under a genuine $q$, so
Theorem~\ref{thm:bound} applies to every summary built this way.

\section{The pair score separates}
\label{sec:separates}

Substituting the merged summary into \eqref{eq:Phi} gives
\begin{equation}\label{eq:sep}
\Phi_k(s_P)=a_{mk}+b_{nk}+\ell_m^\top q_B+\kappa_k(m,n),
\end{equation}
where, with $q_A$ the top distribution of column $m$ and $q_B$ that of
column $n$,
\begin{align*}
a_{mk}&=\Phi_k(s_{A,m})-\tfrac t2\,r_k\,q_A^\top S\pi^k, &
b_{nk}&=\Phi_k(s_{B,n})-\tfrac t2\,r_k\,q_B^\top S\pi^k+\log r_k,\\
\ell_m&=L^\top q_A, &
\kappa_k(m,n)&=-\sum_xq_A(x)\,q_B(x)\log\pi^k_x-(q_A^\top q_B)\log r_k .
\end{align*}
The vectors $a_m,b_n\in\mathbb R^K$ and $\ell_m\in\mathbb R^A$ depend on one
column each and are computed once per column in $O(KA+A^2)$ time, outside
the DP.
With $u_{mk}=w_ke^{a_{mk}}$ and $v_{nk}=e^{b_{nk}}$, stored with a per-column
maximum to avoid underflow, \eqref{eq:mix} becomes
\begin{equation}\label{eq:score}
\log p(D)\ \ge\ \log\sum_ku_{mk}\,v_{nk}\,e^{\kappa_k(m,n)}\;+\;\ell_m^\top q_B .
\end{equation}
Without $\kappa$, this is one length-$K$ and one length-$A$ dot product.

\paragraph{The cross term.} $\kappa_k$ is the only part that couples $m$, $n$
and $k$. It arises from the no-substitution case $x_A=x_B$ of the bridge, in
which $B$'s top token is inherited rather than replaced by a substitution
destination.

\begin{lemma}[sparse cross term]\label{lem:sparse}
Dropping any subset of the terms $-q_A(x)q_B(x)\log\pi^k_x$ from $\kappa_k$
leaves a lower bound on $\log p(D)$.
\end{lemma}
\begin{proof}
Each term is $\ge0$ because $\log\pi^k_x\le0$, and the right-hand side of
\eqref{eq:score} is increasing in every $\kappa_k$.
\end{proof}

So, for a small threshold $\epsilon$ (we use $10^{-3}$), only residues where
both top distributions have non-negligible mass need to be kept, at $K$
operations each. Define the \emph{supports}
$\mathcal S_m=\{x:q_A(x)>\epsilon\}$ and $\mathcal S_n=\{x:q_B(x)>\epsilon\}$,
short lists computed once per column, and keep the terms with
$x\in\mathcal S_m\cap\mathcal S_n$; by the lemma, the result is still a lower
bound. Top distributions are usually peaked, so the intersection is
usually empty or a single
residue; when both tops are point masses on $x$, the factor is a precomputed
$1/\pi^k_x$. The rate part $-(q_A^\top q_B)\log r_k$ can have either sign and
must be kept: one length-$A$ dot product and $K$ exponentials, and nothing at
all with a single rate class.

\paragraph{Cost.} Per DP cell: $O(K+A)$, plus $O(K)$ per residue in the
overlap, against $O(KA)$ for exact pruning. Per column, outside the DP:
$O(K+A)$ storage and $O(KA+A^2)$ precomputation (the $A^2$ is $\ell_m$),
against $KA$ storage and $O(KA^2)$ pre-multiplication for exact pruning. An
alignment of two MSAs of length $L$ has $O(L^2)$ cells but only $O(L)$
columns, so the per-cell costs are the ones that matter.

\section{Re-estimated clade summaries}
\label{sec:reest}

Nothing requires the summaries inside a clade to come from single-jump merges:
Theorem~\ref{thm:bound} holds for any $q$. Once a clade's MSA is fixed, its
column summaries can be recomputed from the exact posterior $p_{k^*}(h\mid D)$
of a reference class $k^*$ (here the clade's best class), rooted at the clade's
top node. The statistics $V'$, $U$, $T$ are the expected sufficient statistics
of that posterior \citep{HolmesRubin2002,HobolthJensen2005}, computed in
eigenspace as in Appendix~\ref{app:hr}; $q_{\mathrm{top}}$
is the posterior at the top node; and $C=\log p_{k^*}(D)-\big(V'\cdot\log
\pi^{k^*}+U\log r_{k^*}-r_{k^*}T\cdot S\pi^{k^*}\big)$, since
$\E_q\log p_{k^*}(h)+H(q)=\log p_{k^*}(D)$ when $q$ is the posterior. The
summary is rooted at the top node rather than at a leaf, so no re-orientation
is needed in Lemma~\ref{lem:merge}: its root token is $e_{x_B}$ itself. This
costs one up--down pass plus per-branch expected statistics per column per
merge, outside the DP; the single-jump construction is then used only for the
final merge, i.e.\ for the DP cell being scored, and its errors no longer
accumulate over merges.

\section{Test on Pfam}
\label{sec:pfam}

We took 60 random families from the
test split of the Pfam seed alignments (54 had at least 6 sequences),
subsampled to at most 64 sequences, and built BioNJ trees \citep{Gascuel1997}
from Kimura protein
distances. The model is LG exchangeabilities with the C20 profiles and
weights, each class normalized to mean rate 1. For 30 random columns per
family, the tree is restricted to the column's non-gap leaves and split at a
random edge into two clades with at least two leaves each. The exact
log-likelihood is computed by per-class pruning.

\begin{center}
\begin{tabular}{lrrrrr}
\toprule
median gap (nats) / leaves & 4--7 & 8--15 & 16--31 & 32--64 & all\\
\midrule
columns & 110 & 250 & 311 & 949 & 1620\\
single-jump summaries throughout & 5.1 & 6.8 & 17.9 & 35.2 & 21.4\\
re-estimated, single-jump final merge & 2.9 & 2.1 & 2.0 & 1.5 & 1.8\\
re-estimated, exact final bridge & 2.8 & 2.0 & 1.9 & 1.4 & 1.7\\
one class's exact posterior, shared & 0.48 & 0.30 & 0.28 & 0.10 & 0.19\\
\bottomrule
\end{tabular}
\end{center}

The rows are: single-jump summaries throughout
(Section~\ref{sec:summaries}); re-estimated clade summaries
(Section~\ref{sec:reest}) with a single-jump or an exact endpoint-conditioned
final bridge; and $q(h)$ equal to the exact posterior of the column's best
class over the whole tree, shared by all classes. The bound was valid on every
column. Four conclusions:
\begin{itemize}
\item \emph{Sharing one $q(h)$ across classes is nearly free} (last row), so
  the class-free structure that gives $O(K+A)$ is not what costs accuracy.
\item \emph{Single-jump summaries accumulate error}, about 0.7 nats per leaf.
  Re-estimating clade summaries (Section~\ref{sec:reest}) removes the
  accumulation, and the gap stays near 1.5--2 nats at every depth.
\item \emph{The single-jump bridge is nearly free}: an exact endpoint-conditioned
  final bridge gains only 0.07 nats, and costs $O(A^2+KA)$ per cell.
\item \emph{The remaining gap is the independence of $x_A$ and $x_B$} at the
  final merge: each clade's top distribution is its posterior given its own
  leaves only. Coupling them is what exact pruning does, and it makes the class
  readout trilinear in (column $m$, column $n$, class), i.e.\ $O(KA)$ per cell.
\end{itemize}

\paragraph{Rankings.} For each column we also scored a misaligned cell, in
which $B$'s residues come from a different random column of the same family,
and compared log-odds match scores $\log p(AB)-\log p(A)-\log p(B)$. For single-jump summaries
the single-clade terms are also single-jump bounds; for re-estimated clades
they are exact, since those clades are pruned anyway to re-estimate them.

\begin{center}
\begin{tabular}{lrrr}
\toprule
 & prefers aligned & sign agrees with exact & Spearman\\
\midrule
exact & 85.4\% & -- & --\\
single-jump throughout & 84.0\% & 91.7\% & 0.910\\
re-estimated clades & 83.3\% & 92.1\% & 0.932\\
\bottomrule
\end{tabular}
\end{center}

\section{Tightness at alignment divergences}
\label{sec:tightness}

The Pfam test measures how far the bound is from the exact log-likelihood of a
column. An aligner needs something else: the \emph{difference} between the
score of matching two columns and the scores of leaving them unmatched must be
right, cell by cell. At the divergences of BAliBASE this fails for the
single-jump construction.

\paragraph{Diagnosis.} For columns of one residue each, the single-jump bridge
makes no substitution when the residues agree and exactly one when they differ,
so it misses multiple substitutions and return paths. Relative to exact pruning, the
log-odds of a match (match score minus the two single-column scores) is biased
low by 0.4, 1.0, 1.7 and 2.9 nats at $t=0.3$, 1, 2 and 4 under C20, and by up to
9 nats for residue pairs with small exchangeability. On the two clades at the
root of the guide tree of 11 BAliBASE families, with columns from the
reference alignment, the bias is 2.7 to 7.5 nats per cell, while its spread
about the bias is only about 1 nat. The single-column scores carry no bridge, so the slack of
the bridge becomes a systematic penalty on every match. Run through the Pair HMM
merge with the MixFrag parameters, the recall of reference column pairs falls
from 0.52 (Viterbi) and 0.55 (expected-SPS decoding,
Section~\ref{sec:decoding}) with exact pruning to 0.02 and 0.23 with the bound.

\paragraph{Design study.} We compared five ways to tighten the scores on this
harness, each implemented, measured, and then checked by independent
reviewers who tried to break the claimed validity and audited the claimed cost;
Section~\ref{sec:choose} compares them with the scores below.
\begin{itemize}
\item \emph{Separable multi-jump bridges} keep the summaries and the $O(K+A)$
  cost: each class uses its own endpoint-conditioned bridge on the final branch
  (a valid mixture bound, since each class's term is a separate ELBO), with a
  separable minorant fitted per branch time and exact residuals on the top-2
  residues. The best bridge family that keeps the class statistics separable
  improves the leaf-pair gap only marginally over the single jump.
\item \emph{Exact bridges on a sparse block} of endpoint pairs, with
  per-column conditional bounds, close much of the gap but cost more per cell
  than exact pruning at $K=A=20$.
\item \emph{Consistent singles} bound the unmatched hypotheses with the same
  kind of slack, so that it cancels in the log-odds; this needs per-class
  partials and a second Forward--Backward pass.
\item A \emph{calibrated} correction of the match log-odds, fitted on Pfam, is
  accurate but is not a bound.
\item The \emph{sparse-class bound} below.
\end{itemize}



\paragraph{The sparse-class bound.} Keep the exact per-class partials $F^k$ of
each column (the $KA$ numbers of exact pruning), and write
$l^k=\pi^k\circ F^k$ and $\mathsf K^k=M^k(t)\,\Pi_k^{-1}$, which is symmetric and
nonnegative by reversibility. The exact class term of a cell is
$E_k=l_A^{k\top}\mathsf K^kl_B^k$. Split each side into its top residue and the
rest, $l_A=v_Ae_{x^*}+a_2$ and $l_B=v_Be_{y^*}+b_2$:
\[
E_k=v_A\,(\mathsf K l_B)(x^*)+v_B\,(\mathsf K l_A)(y^*)-v_Av_B\,\mathsf K(x^*,y^*)
+a_2^\top\mathsf K b_2 ,
\]
where the first three terms are exact and cost $O(1)$ each from per-column
quantities, and
\[
a_2^\top\mathsf Kb_2\;\ge\;\max\big(|a_2|\min_x(\mathsf Kb_2)(x),\;|b_2|\min_y(\mathsf Ka_2)(y)\big)
\]
because every entry is nonnegative. Summing
over classes before taking the maximum of the two remainder bounds keeps a
lower bound and turns the whole tier into five matrix products with $K$
nonzeros per row. Finally, on the union $S$ of the two columns' top two classes
by class posterior, the exact $E_k$ (cost $A$ each) replaces its lower bound.
The result is a genuine lower bound on the exact column-pair likelihood,
exact whenever a column is a single residue, at $5K+|S|(A+O(1))$ operations per
cell ($|S|\approx3.7$ on average) instead of $KA$; the single-column scores are
exact. Its log-odds differ from exact by $0.013$ nats on average (Spearman
correlation $0.9998$), and its recall matches exact pruning, also on a harder
set of 12 families in which both clades have several sequences (0.61 / 0.68
against 0.62 / 0.67). The price is the per-column state: $KA$ numbers instead
of the $O(A)$ summaries.

\paragraph{The posterior-pair score.} The exact match log-odds has a compact form.
Let $\gamma_k$ be a column's class posterior and $q^k$ its class-$k$ posterior
distribution of the top state. Then
\[
\log\frac{p(A,B)}{p(A)\,p(B)}
=\log\sum_k\frac{\gamma^A_k\gamma^B_k}{w_k}\;q^{k\top}_A\,\mathsf R^k\,q^k_B,
\qquad \mathsf R^k=M^k(t)\,\Pi_k^{-1}.
\]
For one class this is $\log q_A^\top\mathsf Rq_B$, a single $A$-dimensional
bilinear form. Collapsing the classes, i.e.\ replacing each column's tops by its
mixture posterior $\bar q=\sum_k\gamma_kq^k$ and $\mathsf R^k$ by the
class-marginal ratio $\bar{\mathsf R}=J(t)\oslash\bar\pi\bar\pi^\top$, gives the
score $\log\bar q_A^\top\bar{\mathsf R}\,\bar q_B$ at $O(A)$ per cell. It is
exact for one class and for single-residue columns, and for $K>1$ it ignores
only that the two columns share a class. It is not a bound, so it serves
decoding, not likelihood computation. On the C20 clade pairs its recall is
0.54 / 0.57 against 0.54 / 0.57 for exact pruning, with a log-odds bias of
$-0.3$ nats and a spread of 0.5; restricting the exact form to a few classes
instead (a valid bound) loses badly (0.28 / 0.49 with each column's top two
classes), because under this normalization the dropped classes get no lower
bound at all.

\paragraph{The log-expectation score.} With the classes collapsed, the
posterior-pair score is
\[
\log\sum_{a,b}\bar q_A(a)\,\bar q_B(b)\,\frac{J_{ab}(t)}{\bar\pi_a\bar\pi_b},
\]
which is the form of MUSCLE's log-expectation profile score
$\log\sum_{ij}f_{xi}f_{yj}\,p_{ij}/(p_ip_j)$ \citep{Edgar2004}, with the column
frequencies $f$ replaced by top-state distributions and the fixed pair
probabilities (from VTML240 in MUSCLE) by the model's time-dependent
class-marginal pair distribution. For one substitution chain and posterior tops
it is exactly the log-odds of matching the two columns, which gives the
log-expectation score a probabilistic derivation; MUSCLE's column-occupancy
factor has its counterpart in the residue-count weights of the expected-SPS
decoder (Section~\ref{sec:decoding}). Nor does the score need the posterior: on
the soft-Fitch tops of single-jump summaries, with no pruning and no per-class
partials, its recall on the C20 clade pairs is 0.53 / 0.59 against 0.54 / 0.57
for exact pruning. In that form the profile mixture enters only through $J(t)$,
and position specificity only through the top distributions. The single-column
scores cancel in decoding (every column contributes one, matched or not), so
only the top distributions are needed. The score does not use the
sufficient-statistic linearization of Section~\ref{sec:bound}, only the soft-Fitch
recursion of Section~\ref{sec:summaries}.

\paragraph{Conclusion.} Bounds built from the $O(A)$ column summaries, with a
bridge whose class statistics are separable, are too loose to drive alignment
decisions at these divergences; the loss comes from the bridge on the final
branch, and the tops being independent, not from sharing one history
distribution across classes. Keeping per-class partials per column and
bounding only the per-cell work gives a lower bound that is practically exact;
collapsing the classes of the exact posterior form gives an approximation that is
as accurate for alignment at $O(A)$ per cell.

\section{Choosing a column-pair score}
\label{sec:choose}

A column-pair score serves three purposes with different requirements.
\emph{Likelihood computation} (parameter estimation, model comparison) needs the
exact value or a genuine lower bound. \emph{Decoding} needs the differences
between cells to be right; an approximation is acceptable, and a bound whose
slack varies across cells is not. \emph{Markov chain Monte Carlo} over
alignments can use any score to propose moves: the Metropolis--Hastings ratio
corrects for the proposal, so validity as a bound is irrelevant, but the
proposal distribution must be close to the target, and the exact score is needed
only along the current and proposed paths, i.e.\ for $O(L)$ cells rather than
$O(L^2)$. We measure closeness by the Kullback--Leibler divergence
$\mathrm{KL}(p\,\|\,q)$ between the Pair HMM posteriors over alignments of a
clade pair under the exact scores ($p$) and under the score being assessed
($q$). Because the two Pair HMMs differ only in their emissions,
$\log p(\text{path})-\log q(\text{path})$ is linear in the path's cell
occupancy, so
\[
\mathrm{KL}(p\,\|\,q)=\sum_{\text{cells}}P_p(\text{cell})\,\Delta(\text{cell})
-\big(\log Z_p-\log Z_q\big),
\]
with $\Delta$ the difference of emission scores and the occupancies
$P_p$ the gradients of $\log Z_p$ (checked against enumeration of all paths of
a small pair). $e^{\mathrm{KL}(p\|q)}$ roughly measures the number of proposals
needed per effective sample of the target; we report it per column, which
indicates the cost of proposing blocks of columns.

\begin{table}[t]
\centering
\footnotesize
\setlength{\tabcolsep}{4pt}
\caption{Column-pair scores on the C20 root clade pairs of 9 BAliBASE families
(re-estimated clade summaries where summaries are used). Operations per cell at
$K=A=20$; state kept per column; recall of reference column pairs by the
MixFrag Pair HMM merge; KL per alignment column between the Pair HMM posteriors
under the exact scores and under the score. For LG ($K=1$) the sparse-class and
posterior-pair scores are exact.}
\label{tab:choose}
\begin{tabular}{lccccc}
\toprule
score & bound & ops/cell & state/column & recall & KL/col\\
\midrule
exact pruning & exact & 400 & $KA$ partials & 0.54 / 0.57 & 0\\
sparse-class bound & yes & $\approx210$ & $KA$ partials & 0.52 / 0.57 & 0.0002\\
posterior-pair (mixture tops) & no & $\approx A$ & $KA$ partials & 0.54 / 0.57 & 0.017\\
log-expectation, posterior tops & no & $\approx A$ & $A$$^\ddagger$ & 0.55 / 0.57 & 0.018\\
log-expectation, soft-Fitch tops$^*$ & no & $\approx A$ & $A$ & 0.53 / 0.59 & 0.020\\
calibrated log-odds & no & $O(K{+}A)$ & summaries & 0.48 / 0.55 & 0.041\\
consistent singles & yes$^\dagger$ & $O(K{+}A)$, 2 passes & $KA$ partials & 0.43 / 0.50 & 0.026\\
separable multi-jump & yes & $\approx80$ & summaries & 0.18 / 0.32 & 0.35\\
single-jump bound & yes & $\approx80$ & summaries & 0.00 / 0.14 & 1.0\\
\bottomrule
\end{tabular}

\smallskip
{\footnotesize Recall is Viterbi / expected-SPS decoding. $^*$On single-jump
summaries, whose tops are soft-Fitch distributions. $^\dagger$Every emission is
a lower bound, but the single-column scores are deliberately loosened; use the
tight ones for likelihoods. $^\ddagger$Plus an exact posterior pass per column
per merge.}
\end{table}

Table~\ref{tab:choose} summarizes the comparison. For likelihoods the
sparse-class bound is the choice: a genuine lower bound, practically exact, at
about half the per-cell cost of exact pruning with the same per-column state.
For decoding, the posterior-pair and log-expectation scores match exact pruning
at $O(A)$ per cell, and the soft-Fitch version needs only $A$ numbers per column
and no pruning at all. For MCMC they are also the natural proposals: about
0.02 nats per column means that a block of 50 columns costs about one nat, at a
twentieth of the per-cell cost of the exact score, with the exact per-class
score evaluated only along the proposed and current paths. The sparse-class
bound is an almost perfect proposal (0.0002 nats per column) but saves only a
factor of two. The bounds built on single-jump summaries, including the
separable multi-jump ones, are poor for all three purposes: hundreds of nats per
clade pair.

\section{Viterbi progressive alignment under MixFrag}
\label{sec:viterbi}

This section assembles the pieces into a complete progressive aligner. The
indel model is MixFrag \citep{HolmesMixFrag}: TKF92 in which each fragment
independently draws a \emph{fragtype} $f\in\{1,\dots,F\}$ with weight $w_f$ and
then has geometric length with extension probability $r_f$; the indel process
and the substitution process are shared across fragtypes, and $F=1$ is TKF92.
The substitution model is the $K$-class profile mixture of
Section~\ref{sec:model}. Everything the pseudocode needs is restated here.

\paragraph{Pair HMM.} As in any progressive aligner, the columns of each
clade's MSA are treated as the sequence at the clade's top node. MixFrag is
time-reversible like TKF92 (its links follow TKF91, and each fragment's
fragtype and length are fixed when it is born), so the joint distribution of the
sequences at sibling nodes $A$ and $B$ is that of a single branch of length
$t=t_A+t_B$, with $A$ playing the ancestor. Let $\lambda<\mu$ be the insertion
and deletion rates, $\kappa=\lambda/\mu$, and
\[
\alpha=e^{-\mu t},\qquad
\beta=\frac{\lambda e^{-\lambda t}-\lambda\alpha}{\mu e^{-\lambda t}-\lambda\alpha},\qquad
\gamma=1-\frac{\mu\beta}{\lambda(1-\alpha)}.
\]
The TKF91 joint Pair HMM has transition probabilities $\tau_{XY}$ over
$X,Y\in\{S,M,I,D,E\}$: the rows $X\in\{S,M,I\}$ are identical,
\[
\tau_{XM}=(1-\beta)\alpha\kappa,\quad
\tau_{XI}=\beta,\quad
\tau_{XD}=(1-\beta)(1-\alpha)\kappa,\quad
\tau_{XE}=(1-\beta)(1-\kappa),
\]
and row $D$ is the same with $\gamma$ in place of $\beta$. The MixFrag Pair HMM
has states $S$, $E$ and $Y_g$ for $Y\in\{M,I,D\}$ and $g\le F$, with
\[
\tau_{S,Y_g}=w_g\tau_{SY},\qquad
\tau_{X_f,Y_g}=r_f\,\delta_{XY}\delta_{fg}+(1-r_f)\,w_g\,\tau_{XY},\qquad
\tau_{X_f,E}=(1-r_f)\,\tau_{XE}.
\]
The self-loop term $r_f\delta_{XY}\delta_{fg}$ extends the current fragment;
every other transition ends it (factor $1-r_f$) and starts a new one of a
freshly drawn fragtype (factor $w_g$). Emissions do not depend on fragtype:
$M$ emits a column of $A$'s MSA together with a column of $B$'s, $D$ a column
of $A$'s alone, and $I$ a column of $B$'s alone.

\paragraph{Emission scores.} The match score of columns $m$ and $n$ is the
bound \eqref{eq:score} on the log-likelihood of the combined column, and the
score of a column on its own is the bound $\log\sum_kw_ke^{\Phi_k(s)}$ of
Theorem~\ref{thm:bound} for its clade. Both are computed from the column
summaries $s=(V',U,T,C,q)$ of Section~\ref{sec:summaries}. The rank-one form of
the fragtype transitions means that Viterbi only has to remember, for each
state type and cell, the best way to \emph{end} a fragment there, so the DP
costs $O(F)$ per state type per cell rather than $O(F^2)$.

\begin{algorithm}[p]
\caption{Per-column preparation for one merge (outside the DP)}
\label{alg:prep}
\begin{algorithmic}[1]
\Require summaries $s_c=(V'_c,U_c,T_c,C_c,q_c)$ for the columns $c$ of one
  clade's MSA; $t=t_A+t_B$; precomputed $\log\pi^k$, $\log r_k$, $\log w_k$,
  $G_{kx}=(S\pi^k)_x$, $L=(\log S+\log t)\circ(1-I)$
\Function{Prepare}{$\{s_c\}$, side}
  \For{each column $c$}
    \State $\Phi_{ck}\gets C_c+V'_c\cdot\log\pi^k+U_c\log r_k-r_k\,T_c\cdot G_k$
      \Comment{all $k$: $O(KA)$}
    \State $\mathrm{single}_c\gets\operatorname{LSE}_k\big(\log w_k+\Phi_{ck}\big)$
      \Comment{emission score of $c$ alone}
    \If{side $=A$}
      \State $a_{ck}\gets\log w_k+\Phi_{ck}-\tfrac t2\,r_k\,q_c\cdot G_k$;
        \quad $\ell_c\gets L^\top q_c$ \Comment{$O(KA+A^2)$}
    \Else
      \State $a_{ck}\gets\Phi_{ck}-\tfrac t2\,r_k\,q_c\cdot G_k+\log r_k$
    \EndIf
    \State $\hat a_c\gets\max_ka_{ck}$;\quad $u_{ck}\gets e^{a_{ck}-\hat a_c}$
      \Comment{scaled to avoid underflow}
    \State $\mathcal S_c\gets\{x:q_c(x)>\epsilon\}$ \Comment{support of the top distribution (Section~\ref{sec:separates})}
  \EndFor
\EndFunction
\end{algorithmic}
\end{algorithm}

\begin{algorithm}[p]
\caption{Match score of column $m$ of $A$ and column $n$ of $B$ (inside the DP)}
\label{alg:pair}
\begin{algorithmic}[1]
\Function{PairScore}{$m,n$}
  \State $\sigma\gets q_m\cdot q_n$ \Comment{$O(A)$}
  \State $\chi_k\gets-\sigma\log r_k$ for all $k$ \Comment{rate cross term; zero if all $r_k=1$}
  \For{$x\in\mathcal S_m\cap\mathcal S_n$} \Comment{sparse $\log\pi$ cross term; usually 0--1 residues}
    \State $\chi_k\gets\chi_k-q_m(x)\,q_n(x)\log\pi^k_x$ for all $k$
  \EndFor
  \State \Return $\hat a_m+\hat b_n+\log\sum_ku_{mk}\,v_{nk}\,e^{\chi_k}+\ell_m\cdot q_n$
    \Comment{$O(K+A)$}
\EndFunction
\end{algorithmic}
\end{algorithm}

Here $u$, $\hat a$ are the \textsc{Prepare} outputs for $A$'s columns and $v$,
$\hat b$ those for $B$'s, and $\mathcal S_m$, $\mathcal S_n$ the supports of
the two columns' top distributions, also from \textsc{Prepare}. By
Lemma~\ref{lem:sparse}, restricting the $\log\pi$ cross term to
$\mathcal S_m\cap\mathcal S_n$ keeps the score a lower bound.

\begin{algorithm}[p]
\caption{Viterbi alignment of two clade MSAs under the MixFrag Pair HMM}
\label{alg:viterbi}
\begin{algorithmic}[1]
\Require $L_A,L_B$ columns; emission scores $e_M(i,j)=\textsc{PairScore}(i,j)$,
  $e_D(i)=\mathrm{single}^A_i$, $e_I(j)=\mathrm{single}^B_j$; all logs of
  $\tau$, $r_f$, $1-r_f$, $w_g$
\Function{Viterbi}{}
  \State predecessor offsets: $\Delta_M=(1,1)$, $\Delta_D=(1,0)$, $\Delta_I=(0,1)$
  \State $H_S(0,0)\gets0$; all other $H_X$, $V_{Y,g}$ $\gets-\infty$
    \Comment{$H_X(i,j)$: best score ending a fragment in state $X$ at $(i,j)$}
  \For{$i=0,\dots,L_A$}
    \For{$j=0,\dots,L_B$}
      \For{$Y\in\{M,D,I\}$ with $(i',j')=(i,j)-\Delta_Y\ge(0,0)$}
        \State $\mathrm{new}\gets\max_{X\in\{S,M,I,D\}}\big(H_X(i',j')+\log\tau_{XY}\big)$;
          \ record the argmax $X^\star_Y(i,j)$
        \For{$g=1,\dots,F$}
          \State $\mathrm{ext}\gets V_{Y,g}(i',j')+\log r_g$;
            \quad $\mathrm{fresh}\gets\mathrm{new}+\log w_g$
          \State $V_{Y,g}(i,j)\gets e_Y(i,j)+\max(\mathrm{ext},\mathrm{fresh})$;
            \ record which of the two won
        \EndFor
        \State $H_Y(i,j)\gets\max_g\big(V_{Y,g}(i,j)+\log(1-r_g)\big)$;
          \ record the argmax $g^\star_Y(i,j)$
      \EndFor
    \EndFor
  \EndFor
  \State $\mathrm{best}\gets\max_{X\in\{M,I,D\}}\big(H_X(L_A,L_B)+\log\tau_{XE}\big)$
    \Comment{$L_A,L_B\ge1$, so $S\to E$ is impossible}
  \State \textbf{traceback:} start in $(X,g^\star_X(L_A,L_B))$ at $(L_A,L_B)$ for the
    maximizing $X$. From $(Y,g)$ at $(i,j)$, output the column type $Y$ and move to
    $(i',j')=(i,j)-\Delta_Y$; if ``ext'' won, stay in $(Y,g)$; otherwise go to
    $X=X^\star_Y(i,j)$ in fragtype $g^\star_X(i',j')$, and stop on reaching $S$.
  \State \Return best, and the column sequence read in reverse
\EndFunction
\end{algorithmic}
\end{algorithm}

The DP makes $3F$ updates and three emission evaluations per cell: $O(F+K+A)$
per cell, $O(L_AL_B(F+K+A))$ per merge. Only the $H_X$ values of the
predecessor row are needed for the scores, but the back-pointers
($X^\star_Y$, $g^\star_Y$ and one bit per $(Y,g)$) take $O(L_AL_BF)$ memory.

\begin{algorithm}[p]
\caption{Summaries of the parent MSA's columns}
\label{alg:newsum}
\begin{algorithmic}[1]
\Function{Merge}{$s_A,s_B,t_A,t_B$} \Comment{single-jump merge, Section~\ref{sec:summaries}}
  \State $t\gets t_A+t_B$;\quad $q_A,q_B\gets$ the two top distributions
  \State $V'\gets V'_A+V'_B-q_A\circ q_B$;\quad $U\gets U_A+U_B+1-q_A\cdot q_B$
  \State $T\gets T_A+T_B+\tfrac t2(q_A+q_B)$;\quad $C\gets C_A+C_B+q_A^\top Lq_B$
  \State $q\gets q_A\circ q_B+\tfrac{t_B}{t}q_A\circ(1-q_B)+\tfrac{t_A}{t}q_B\circ(1-q_A)$
  \State \Return $(V',U,T,C,q)$
\EndFunction
\Function{Extend}{$s,t_A$} \Comment{column present on one side only}
  \State \Return $s$ with $T\gets T+t_A\,q$ \Comment{no substitution on the branch up to the parent; top unchanged}
\EndFunction
\Function{Reestimate}{column $c$, its tree rooted at the parent $P$, cheap summary $s$}
  \State $k^\star\gets\arg\max_k\big(\log w_k+\Phi_k(s)\big)$ \Comment{class chosen from the cheap summary}
  \State $(q,V',U,T,\log p_{k^\star}(c))\gets\textsc{CladeStatistics}(c,\ \text{tree rooted at }P,\ k^\star)$
    \Comment{Algorithm~\ref{alg:clade}}
  \State $C\gets\log p_{k^\star}(c)-\big(V'\cdot\log\pi^{k^\star}+U\log r_{k^\star}
    -r_{k^\star}T\cdot G_{k^\star}\big)$
  \State \Return $(V',U,T,C,q)$ \Comment{$O(n_cA^2+A^3)$ for $n_c$ leaves}
\EndFunction
\end{algorithmic}
\end{algorithm}

\begin{algorithm}[p]
\caption{Progressive alignment}
\label{alg:progressive}
\begin{algorithmic}[1]
\Require rooted binary guide tree with branch lengths (Section~\ref{sec:distances});
  sequences at its leaves;
  MixFrag parameters $(\lambda,\mu,\{r_f,w_f\})$; profile mixture $(S,\{w_k,\pi^k,r_k\})$
\For{each leaf $\ell$}
  \State MSA$_\ell\gets$ the sequence itself; each residue $x$ gets the leaf summary
    $(e_x,0,0,0,e_x)$
\EndFor
\For{each internal node $P$ in postorder, with children $A,B$ at distances $t_A,t_B$}
  \State $\textsc{Prepare}(\text{columns of }A,A)$;\quad $\textsc{Prepare}(\text{columns of }B,B)$
  \State (best, path) $\gets\textsc{Viterbi}()$
  \For{each step of the path}
    \If{$M(m,n)$}
      \State new column = rows of $A$'s column $m$ and $B$'s column $n$;
        \ $s\gets\textsc{Merge}(s_{A,m},s_{B,n},t_A,t_B)$
    \ElsIf{$D(m)$}
      \State new column = $A$'s column $m$ with gaps in $B$'s rows;
        \ $s\gets\textsc{Extend}(s_{A,m},t_A)$
    \Else \Comment{$I(n)$}
      \State new column = $B$'s column $n$ with gaps in $A$'s rows;
        \ $s\gets\textsc{Extend}(s_{B,n},t_B)$
    \EndIf
    \State optionally $s\gets\textsc{Reestimate}(\text{new column},\ \text{tree below }P,\ s)$
      \Comment{removes accumulated error}
  \EndFor
\EndFor
\State \Return the MSA at the root
\end{algorithmic}
\end{algorithm}

\paragraph{Remarks.}
\begin{itemize}
\item \textsc{Extend} is a valid choice of $q$ (no substitution on the branch
  from the child to $P$), so every summary stays exact bookkeeping for a
  genuine $q$ and Theorem~\ref{thm:bound} applies throughout.
\item Without \textsc{Reestimate}, the per-merge cost outside the DP is
  $O((L_A+L_B)(KA+A^2))$. With it, add $O(n_cA^2+A^3)$ per column of the parent
  MSA; the Pfam test (Section~\ref{sec:pfam}) shows this removes error that
  otherwise accumulates at about 0.7 nats per leaf.
\item With $F=1$ ($w_1=1$, $r_1=r$) MixFrag is TKF92, and
  Algorithm~\ref{alg:viterbi} becomes the standard three-state TKF92 Viterbi:
  $H_X=V_X+\log(1-r)$, and each state either extends its fragment
  ($\log r$) or starts a new one ($\log\tau_{XY}$ from the fragment-end score).
  With $K=1$ as well, the emissions are single-GTR bounds, i.e.\ a TKF92
  profile aligner without position-specific scoring.
\item Replacing $\max$ by log-sum-exp in Algorithm~\ref{alg:viterbi} gives the
  Forward algorithm, which sums over alignments and fragtypes, and
  posterior-decoding variants follow in the usual way.
\end{itemize}

\section{Decoding objectives}
\label{sec:decoding}

Algorithm~\ref{alg:viterbi} maximizes over complete state paths, which fixes
the fragtype of every column as well as the order of insertions and deletions
within each gap. Neither is part of an alignment. We compared the natural
alternatives on the MixFrag Pair HMM with the Pfam-trained parameters, by exact
enumeration of all state paths on 9{,}000 short pairs and on 540 pairs of
realistic length; every exact algorithm below agreed with enumeration to
$10^{-14}$.

\paragraph{Objectives.} Let $\mathrm{O1}$ be the Viterbi objective (maximum over
state paths); $\mathrm{O2}$ the maximum over column-type sequences (alignments
including the order of gap columns) of the probability summed over fragtype
labels; and $\mathrm{O3}$ the maximum over homology skeletons (the set of matched
pairs) of their probability, also summed over the orders of insertions and
deletions between consecutive matches. Always
$\mathrm{O1}\le\mathrm{O2}\le\mathrm{O3}\le\log p(x,y)$.

\paragraph{Fragtypes.} Because every transition between column types is rank
one in MixFrag (the new fragment's type is redrawn from $w$), the sum over
fragtype labels of an alignment factorizes over its maximal runs of one column
type, with run factor $\varphi_a(k)=w^\top\mathsf A_a^{k-1}(1-r)$. Hence
$\mathrm{O2}$ is computed exactly by an explicit-duration (semi-Markov) Viterbi
over runs. With $s_f=\sqrt{w_f(1-r_f)}$, $\varphi_a(k)=s^\top(\mathrm{diag}(r)
+\tau_{aa}ss^\top)^{k-1}s$ is a positive mixture of $F$ geometric terms, so run
lengths can be capped at about 30 with a geometric tail at a cost of
$10^{-4}$ nats per run. For labelled HMMs in general this problem is NP-hard
\citep{LyngsoPedersen2002,Brejova2007}; MixFrag escapes because of the rank-one
transitions. Viterbi's maximum over fragtypes is systematically unfair to
MixFrag: with the Pfam parameters it scores every alignment below
$\mathrm{O2}$ by about 0.1 nats per run of length one, 1.5 nats per match run of
length 10 and 9--11 nats per match run of length 100, whereas for TKF92
$\mathrm{O1}=\mathrm{O2}$.

\paragraph{Gap orderings.} $\mathrm{O3}$ is computed exactly by a DP over
alternating match runs and gap blocks, in which a gap block of $a$ deletions and
$b$ insertions contributes the order-summed gap factor $G(a,b)$ of the
summarized-count likelihood, at $O(L^4)$ time. $\log G$ is nearly separable
(residual at most 0.02 nats for $a,b\le40$, $t\le2.5$), so an $\mathrm{O2}$-cost
variant that folds both single-alternation orders into a canonical order is
within $\sum\log(1+\delta)\le0.065$ nats per mixed gap of $\mathrm{O3}$ at
$t\le1$. The order sum gains 0.4--0.6 nats per mixed gap over the best single
order.

\paragraph{The naive hybrid.} Taking the log-sum-exp over predecessors at insert
and delete states and the maximum at match states is a common shortcut. It
computes the best \emph{policy bundle}: one predecessor fixed at each match
state and all predecessors kept at each gap cell. It satisfies
$\mathrm{O1}\le H\le\log p(x,y)$ but is neither above nor below $\mathrm{O2}$ or
$\mathrm{O3}$ (strict inequalities in both directions in thousands of the
brute-force instances, even for TKF92): a gap cell's sum pools paths that entered
the gap from different match cells, i.e.\ different skeletons, while a match
cell keeps only one fragtype and one last gap state. It has no well-defined
traceback through gap cells, and its greedy traceback gives a skeleton that is
not $\mathrm{O3}$-optimal in 4\% of short pairs and 8--15\% of realistic ones.
It maximizes no clean objective.

\paragraph{Posterior decoding.} Summing everything gives the posterior
probability $p_{mn}$ that columns $m$ and $n$ are matched, which is the gradient
of the log Forward probability with respect to the match table, computed by
reverse-mode differentiation of the checkpointed wavefront Forward pass.
Matching the two columns creates $n_mn_n$ new residue pairs (the numbers of
residues in each column), each correct with probability about $p_{mn}$, so the
merge that maximizes the expected sum-of-pairs score maximizes
$\sum_{\text{matched}}p_{mn}n_mn_n$, a Needleman--Wunsch problem with no gap
penalty that the same wavefront kernel solves in the max-plus semiring. This is
expected-SPS decoding at each sibling merge: a form of maximum expected accuracy
decoding \citep{HolmesDurbin1998,Do2005} in which the accuracy is the sum-of-pairs
score, so that, unlike alignment metric accuracy, correctly unaligned residues
earn nothing and there is no gap term. It is greedy (each merge maximizes its own
expected contribution given the clade alignments below it), and it sums over
fragtypes, gap orderings and skeletons, which is what the objectives above
approximate.

\section{Pairwise distances and the guide tree}
\label{sec:distances}

Algorithm~\ref{alg:progressive} needs a rooted guide tree with branch
lengths. We estimate the evolutionary time $t$ separating every pair of
input sequences under the same Pair HMM, build a tree from these distances
with BioNJ \citep{Gascuel1997}, and root it at the midpoint of its longest
path. The distance estimate is one iteration of expectation--maximization
(EM) for $t$, in the style used by FSA \citep{Bradley2009}: a single
Forward--Backward pass at a fixed initial time, followed by maximization of
the expected complete-data log-likelihood, in which the complete data are
the alignment and the expected counts are held fixed.

\paragraph{Model.} Distances use TKF92 (MixFrag with one fragtype) and a
single substitution chain, LG, whatever column model the aligner itself uses:
this is the model under which FSA estimates distances, and it keeps the guide
tree independent of the profile mixture. With $\pi$ and $M(t)$ the LG
equilibrium and transition matrix, the match emissions are
$J_{ab}(t)=\pi_aM_{ab}(t)$ and the insert and delete emissions $\pi_x$.

\paragraph{E-step.} Fix $t_0$ (we use $t_0=1$; Appendix~\ref{app:kgram}
gives an alignment-free alternative). For a pair of sequences $X$, $Y$, run
Forward--Backward on the TKF92 Pair HMM at $t_0$. This gives the expected number
$n_{XY}$ of transitions from state $X$ to state $Y$ and the expected number
$W_{ab}$ of aligned pairs $(a,b)$. Both are gradients of the log Forward
probability, with respect to $\log\tau_{XY}$ and to $\log J_{ab}$
respectively, so reverse-mode differentiation of the linear-space Forward
pass computes them without storing per-cell posteriors.

\paragraph{M-step.} With the counts fixed, the expected complete-data
log-likelihood as a function of $t$ is
\begin{equation}\label{eq:qt}
\mathcal Q(t)=\sum_{X,Y}n_{XY}\log\tau_{XY}(t)+\sum_{a,b}W_{ab}\log J_{ab}(t)
+\text{const},
\end{equation}
where the constant collects the insert and delete emissions, which do not
depend on $t$. We maximize $\mathcal Q$ over $u=\log t$ by a fixed number of
Newton--Raphson steps, with the step clipped and $t$ kept in a bounded
interval. Derivatives come from automatic differentiation, and each step
costs $O(A^3)$ for $M(t)$, independent of the sequence lengths.

\paragraph{Parallelism.} Every pair is independent. Pairs are grouped by
their lengths rounded up to a geometric grid, padded to the grid values, and
processed in batches, with the true lengths passed as traced values so that
one compiled kernel serves a whole batch; the Forward pass runs over
anti-diagonals, whose cells are independent. The Newton steps act on the
fixed-size counts $(n,W)$ and are vectorized over all pairs at once.

\begin{algorithm}[tbp]
\caption{Pairwise distance by one EM iteration}
\label{alg:distance}
\begin{algorithmic}[1]
\Require sequences $X,Y$; TKF92 parameters; LG; $t_0$;
  number of Newton steps $n_{\mathrm{N}}$ (we use 5)
\Function{PairDistance}{$X,Y$}
  \State $\ell(\theta)\gets$ log Forward probability of $(X,Y)$ at $t_0$, as a
    function of $\theta=(\log\tau(t_0),\log J(t_0))$
    \Comment{anti-diagonal, linear space}
  \State $(n,W)\gets\nabla_\theta\,\ell$ \Comment{expected transition and
    aligned-pair counts}
  \State $u\gets\log t_0$
  \For{$1,\dots,n_{\mathrm{N}}$}
    \State $g\gets\mathcal Q'(u)$;\quad $h\gets\mathcal Q''(u)$
      \Comment{Equation~\eqref{eq:qt} as a function of $u=\log t$}
    \State $u\gets u-\mathrm{clip}(g/h,\,-1,\,1)$ if $h<0$, else
      $u\gets u+\mathrm{clip}(g,\,-1,\,1)$
    \State $u\gets\mathrm{clip}(u,\ \log t_{\min},\ \log t_{\max})$
  \EndFor
  \State \Return $e^u$
\EndFunction
\end{algorithmic}
\end{algorithm}

\section{Results on BAliBASE}
\label{sec:results}

We implemented the aligner in JAX: anti-diagonal wavefront Forward and Viterbi
kernels with geometric length padding, pairwise distances by
Algorithm~\ref{alg:distance} batched by padded length, BioNJ guide trees, and
the scorings and decoders above. The indel parameters are the Pfam-trained
TKF92 and MixFrag ($\mathcal F=2$) fits of paper~1; LG is used alone or with
the C20 profiles, the latter with a substitution clock scaled by the family's
C20/LG time ratio (the indel parameters were fitted on the LG clock). We scored
the 120 families of BAliBASE~3 \texttt{bali3pdbm} by sum-of-pairs (SP) and
total-column (TC) scores on the reference core columns, and ran MUSCLE~5.3 and
MAFFT~7.526 through the same scorer. MUSCLE reproduces paper~1's numbers
exactly; so does MAFFT, but only when run with \texttt{--auto}, whereas papers~1
and~3 label the row ``FFT-NS-2''.

\begin{table}[t]
\centering
\footnotesize
\setlength{\tabcolsep}{4pt}
\caption{Accuracy on BAliBASE~3 \texttt{bali3pdbm}: all 120 families, and the
22 families whose sequences are all shorter than 150 residues (the subset of
paper~3). ESPS: expected-SPS decoding; refinement: two rounds of
Section~\ref{sec:refine}. FSA rows are quoted from papers~1 and~3 (sequence
annealing of pairwise posteriors).}
\label{tab:results}
\begin{tabular}{llcccc}
\toprule
& & \multicolumn{2}{c}{120 families} & \multicolumn{2}{c}{22 families, $L<150$}\\
\cmidrule(lr){3-4}\cmidrule(lr){5-6}
method & scoring & SP & TC & SP & TC\\
\midrule
MUSCLE 5.3 (default) & & \textbf{0.841} & \textbf{0.715} & \textbf{0.760} & \textbf{0.662}\\
MAFFT (\texttt{--auto}) & & 0.797 & 0.653 & 0.691 & 0.558\\
MAFFT (FFT-NS-2) & & 0.760 & 0.595 & 0.731 & 0.613\\
\midrule
FSA, TKF92 (paper 1 / 3) & & 0.776 & 0.622 & 0.724 & 0.631\\
FSA, MixFrag $\mathcal F{=}2$ (paper 1) & & 0.790 & 0.640 & -- & --\\
FSA, TKF92-$K{=}20$ (paper 3) & & -- & -- & 0.732 & 0.637\\
\midrule
TKF92, LG, Viterbi & exact & 0.735 & 0.566 & 0.709 & 0.606\\
TKF92, LG, ESPS & exact & 0.748 & 0.582 & 0.734 & 0.638\\
MixFrag, LG, Viterbi & exact & 0.665 & 0.478 & 0.600 & 0.442\\
MixFrag, LG, ESPS & exact & 0.751 & 0.587 & 0.726 & 0.633\\
MixFrag, LG, ESPS, refinement & exact & 0.738 & 0.569 & 0.650 & 0.515\\
MixFrag, LG, ESPS & log-expectation & 0.718 & 0.554 & 0.707 & 0.616\\
\midrule
TKF92, C20, ESPS & exact & 0.730 & 0.562 & 0.699 & 0.592\\
MixFrag, C20, Viterbi & exact & 0.630 & 0.451 & 0.569 & 0.444\\
MixFrag, C20, ESPS & exact & 0.742 & 0.569 & 0.703 & 0.604\\
MixFrag, C20, Viterbi & sparse-class bound & 0.608 & 0.435 & 0.553 & 0.424\\
MixFrag, C20, ESPS & sparse-class bound & 0.736 & 0.567 & 0.702 & 0.602\\
MixFrag, C20, Viterbi & posterior-pair & 0.655 & 0.471 & 0.601 & 0.447\\
MixFrag, C20, ESPS & posterior-pair & 0.744 & 0.583 & 0.696 & 0.586\\
MixFrag, C20, ESPS, refinement & posterior-pair & 0.735 & 0.560 & 0.687 & 0.562\\
MixFrag, C20, ESPS & log-expectation & 0.736 & 0.573 & 0.695 & 0.576\\
MixFrag, C20, ESPS & single-jump bound & 0.244 & 0.121 & 0.277 & 0.137\\
\bottomrule
\end{tabular}
\end{table}

Table~\ref{tab:results} gives the results. Five findings:
\begin{itemize}
\item \emph{Scoring.} Every score that is close to exact on the clade pairs
  (exact pruning, the sparse-class bound, the posterior-pair and
  log-expectation scores) gives SP within 0.015 of the others; the $O(A)$
  scores lose essentially nothing. The single-jump bound fails (SP 0.24), as
  the harness predicted.
\item \emph{Decoding.} Expected-SPS decoding beats Viterbi in every
  configuration. Viterbi is particularly bad for MixFrag (0.665 against 0.735
  for TKF92 with LG), as Section~\ref{sec:decoding} predicts from its maximum
  over fragtypes; with expected-SPS decoding MixFrag is the best of our
  configurations, 0.016 SP above TKF92 Viterbi (better in 72 families, worse
  in 33; Wilcoxon $p=0.001$).
\item \emph{Profiles.} The C20 mixture does not improve on LG here (0.742
  against 0.751 with exact scoring), despite its better fit to single columns.
\item \emph{Refinement.} Two rounds of refinement lower SP (0.751 to 0.738;
  0.726 to 0.650 on the short families), although every step increases the
  model's expected SPS: the posteriors are not calibrated well enough against
  the structural references for coordinate ascent on them to pay.
\item \emph{Comparisons.} The best configuration is 0.09 SP below MUSCLE on the
  full set (worse in 94 families of 120) and below FSA's sequence annealing of
  pairwise posteriors under the same indel models (0.790 for MixFrag). On the
  22 short families it is level with FSA (TKF92: 0.734 against 0.724) and
  0.026 SP below MUSCLE.
\end{itemize}
Our aligner takes 40--100\,s per family on a shared CPU, most of it JAX
compilation per padded shape; the wavefront kernels are written for GPUs,
where each anti-diagonal is one vectorized step.

\paragraph{Refinement.}
\label{sec:refine}
After the progressive pass, each edge of the guide tree splits the sequences in
two; the two sub-alignments (restricted to their rows, all-gap columns removed)
are re-pruned on their sides of the edge, rooted at its endpoints, and realigned
by expected-SPS decoding at that edge's branch length. The old alignment of the
two halves is a feasible path, so each step can only increase the expected
number of correct cross pairs under that step's posteriors; a round visits every
edge.

\section{Open questions}
\begin{itemize}
\item \emph{The gap to MUSCLE and FSA.} Expected-SPS decoding is greedy, one
  merge at a time; FSA's sequence annealing of pairwise posteriors, under the
  same indel models, is 0.04 SP better on the full set. Consistency
  transformations or annealing over the merged profiles' posteriors are the
  natural next steps.
\item \emph{Calibration.} Refinement increases the model's expected SPS but
  lowers the true SPS, so the posteriors are overconfident or biased against
  the structural references. Calibrating them (e.g.\ tempering, or an indel
  model fitted to alignment accuracy rather than likelihood) would help both
  refinement and decoding.
\item \emph{Why C20 does not help.} The profile mixture fits single columns
  better but does not improve alignment accuracy here; whether this is the
  clock (the indel parameters were fitted with LG), the per-class rate
  normalization, or the absence of rate heterogeneity is untested.
\item \emph{MCMC.} The posterior-pair and log-expectation scores are cheap,
  accurate proposals (Section~\ref{sec:choose}); a block-proposal sampler over
  alignments, with exact per-class scores along proposed paths, is untested.
\item \emph{Hardware.} The wavefront kernels are designed for GPUs; the
  benchmark ran on a shared CPU.
\end{itemize}

\appendix
\section{Expected substitution statistics in eigenspace}
\label{app:hr}

This appendix gives the computation behind \textsc{Reestimate}: the expected
dwell times and substitution counts of the exact posterior of one class on one
column \citep{HolmesRubin2002,HobolthJensen2005}. Following
\citet{HolmesRubin2002}, the counts are accumulated over
branches in the eigenbasis of the rate matrix, as ``eigensubstitutions'', and
transformed back to residues once per column. Each branch then costs $O(A^2)$
and only the final back-transformation costs $O(A^3)$.

Drop the class index: $Q_{xy}=rS_{xy}\pi_y$ ($x\ne y$), $Q_{xx}=-r(S\pi)_x$.

\paragraph{Symmetric eigendecomposition.} With $\Pi=\mathrm{diag}(\pi)$,
$B=\Pi^{1/2}Q\,\Pi^{-1/2}$ is symmetric:
$B_{xy}=rS_{xy}\sqrt{\pi_x\pi_y}$ for $x\ne y$ and $B_{xx}=Q_{xx}$. Write
$B=\Psi\Lambda\Psi^\top$ with $\Psi$ orthogonal and $\Lambda=\mathrm{diag}(\lambda_i)$.
Then $M(t)=e^{tQ}=\Pi^{-1/2}\Psi e^{t\Lambda}\Psi^\top\Pi^{1/2}$, so a
matrix--vector product with $M(t)$ or $M(t)^\top$ costs $O(A^2)$ through the
eigenbasis, and no $A\times A$ transition matrix need ever be formed.

\paragraph{Messages.} Root the column's tree at $P$. For a node $v$ with
parent $u$, on a branch of length $t_v$, let $L_v$ be the inside vector
(likelihood of the leaves below $v$ given the state at $v$), $\eta_u$ the
outside vector of $u$ (joint probability of the state at $u$ and all leaves
not below $u$; $\eta_P=\pi$), and $R_v$ the joint probability of the state at
$u$ and all leaves not below $v$, i.e.\ $\eta_u$ times the inside messages
from $u$'s other children. Then the posterior probability that the branch starts
in $a$ and ends in $b$ is $P_{ab}=R_{v,a}M_{ab}(t_v)L_{v,b}/Z_v$, where
$Z_v=R_v^\top M(t_v)L_v=p(D)$ (up to the scaling of $R_v$ and $L_v$). Define
the projections
\[
\tilde R_v=\Psi^\top\Pi^{-1/2}R_v,\qquad \tilde L_v=\Psi^\top\Pi^{1/2}L_v,
\qquad Z_v=\sum_i\tilde R_{v,i}\,e^{\lambda_it_v}\,\tilde L_{v,i}.
\]

\begin{lemma}[eigensubstitution counts]\label{lem:hr}
Let
\[
\mathcal C=\sum_{\text{branches }v}\frac{1}{Z_v}
\big(\tilde R_v\tilde L_v^\top\big)\circ\mathcal D(t_v),
\qquad
\mathcal D_{ij}(t)=\int_0^te^{\lambda_iu}e^{\lambda_j(t-u)}\,du
=\begin{cases}\dfrac{e^{\lambda_it}-e^{\lambda_jt}}{\lambda_i-\lambda_j}&\lambda_i\ne\lambda_j,\\[6pt]
te^{\lambda_it}&\lambda_i=\lambda_j,\end{cases}
\]
and $\tilde J=\Psi\,\mathcal C\,\Psi^\top$. Then the posterior expected total time
in state $x$, summed over branches, is $\tilde J_{xx}$, and the posterior
expected total number of $x\to y$ substitutions is $B_{xy}\tilde J_{xy}$.
\end{lemma}
\begin{proof}
Write $I^{ab}_{xy}(t)=\int_0^tM_{ax}(u)M_{yb}(t-u)\,du$. On one branch,
conditioned on endpoints $(a,b)$, the expected time in $x$ is
$I^{ab}_{xx}/M_{ab}$ and the expected number of $x\to y$ jumps is
$Q_{xy}I^{ab}_{xy}/M_{ab}$ \citep{HolmesRubin2002,HobolthJensen2005}.
Averaging over $P_{ab}$, both are read off
\[
J_{xy}=\sum_{a,b}\frac{P_{ab}}{M_{ab}}I^{ab}_{xy}
=\frac1{Z_v}\sum_{a,b}R_{v,a}\,L_{v,b}\,I^{ab}_{xy},
\]
since $P_{ab}/M_{ab}=R_{v,a}L_{v,b}/Z_v$ has rank one. Substituting
$M_{ax}(u)=\sqrt{\pi_x/\pi_a}\sum_i\Psi_{ai}\Psi_{xi}e^{\lambda_iu}$ and
$M_{yb}(t-u)=\sqrt{\pi_b/\pi_y}\sum_j\Psi_{yj}\Psi_{bj}e^{\lambda_j(t-u)}$,
\[
J_{xy}=\sqrt{\pi_x/\pi_y}\;\frac1{Z_v}\sum_{i,j}\Psi_{xi}\,\tilde R_{v,i}\,
\mathcal D_{ij}(t_v)\,\tilde L_{v,j}\,\Psi_{yj}.
\]
Summing over branches gives $J_{xy}=\sqrt{\pi_x/\pi_y}\,\tilde J_{xy}$. So the
total time in $x$ is $J_{xx}=\tilde J_{xx}$, and the total number of $x\to y$
jumps is $Q_{xy}J_{xy}=rS_{xy}\pi_y\sqrt{\pi_x/\pi_y}\,\tilde J_{xy}
=B_{xy}\tilde J_{xy}$.
\end{proof}

$\tilde R_v$ and $\tilde L_v$ are the ``left and right eigenbasis projections''
of the up and down vectors of \citet{HolmesRubin2002}, and $\mathcal C$ is
their matrix of eigenbasis counts. Each branch adds a rank-one matrix masked
by $\mathcal D(t_v)$, at $O(A^2)$; the back-transformation $\Psi\mathcal C\Psi^\top$
is done once. In floating point, use the second case of $\mathcal D$ when
$|\lambda_i-\lambda_j|\,t$ is below a small tolerance.

\begin{algorithm}[tbp]
\caption{Posterior statistics of one column under one class, in eigenspace}
\label{alg:clade}
\begin{algorithmic}[1]
\Require column $c$; tree restricted to $c$'s non-gap leaves, rooted at $P$
  (which may have one child); class $k$ with $(\pi,r)$ and eigendecomposition
  $B=\Psi\Lambda\Psi^\top$
\Function{CladeStatistics}{$c$, tree, $k$}
  \For{each node $v$ in postorder} \Comment{inside pass}
    \State $L_v\gets e_{c_v}$ if $v$ is a leaf, else $\mathbf 1$;\quad $\sigma_v\gets0$
    \For{each child $w$ of $v$}
      \State $m_w\gets\Pi^{-1/2}\Psi\big(e^{t_w\Lambda}\tilde L_w\big)$
        \Comment{$=M(t_w)L_w$, with $\tilde L_w$ from line 8}
      \State $L_v\gets L_v\circ m_w$;\quad $\sigma_v\gets\sigma_v+\sigma_w$
    \EndFor
    \State $z\gets\max_xL_v(x)$;\quad $L_v\gets L_v/z$;\quad $\sigma_v\gets\sigma_v+\log z$
    \State $\tilde L_v\gets \Psi^\top\Pi^{1/2}L_v$ \Comment{right eigenbasis projection}
  \EndFor
  \State $\log p\gets\log(\pi\cdot L_P)+\sigma_P$;\quad
    $q\gets\pi\circ L_P/(\pi\cdot L_P)$ \Comment{posterior at $P$}
  \State $\mathcal C\gets0$ ($A\times A$);\quad $\eta_P\gets\pi$
  \For{each node $u$ in preorder, and each child $v$ of $u$}
    \Comment{outside pass}
    \State $R_v\gets \eta_u\circ\prod_{v'\ne v}m_{v'}$
      \Comment{over the other children $v'$ of $u$}
    \State $\tilde R_v\gets \Psi^\top\Pi^{-1/2}R_v$ \Comment{left eigenbasis projection}
    \State $Z_v\gets\sum_i\tilde R_{v,i}\,e^{\lambda_it_v}\,\tilde L_{v,i}$
    \State $\mathcal C\gets\mathcal C+\big(\tilde R_v\tilde L_v^\top\big)\circ\mathcal D(t_v)/Z_v$
      \Comment{eigensubstitutions: $O(A^2)$}
    \State $\eta_v\gets\Pi^{1/2}\Psi\big(e^{t_v\Lambda}\tilde R_v\big)$, normalized to sum to 1
      \Comment{$=M(t_v)^\top R_v$}
  \EndFor
  \State $\tilde J\gets \Psi\,\mathcal C\,\Psi^\top$ \Comment{back-transformation, once: $O(A^3)$}
  \State $T_x\gets\tilde J_{xx}$;\quad $N_{xy}\gets B_{xy}\tilde J_{xy}$ ($x\ne y$)
  \State $V'_y\gets q_y+\sum_{x\ne y}N_{xy}$;\quad $U\gets\sum_{x\ne y}N_{xy}$
  \State \Return $(q,V',U,T,\log p)$
\EndFunction
\end{algorithmic}
\end{algorithm}

Every per-branch step is a matrix--vector product through the eigenbasis or a
rank-one update, so \textsc{CladeStatistics} costs $O(n_cA^2+A^3)$ for a
column with $n_c$ leaves: $O(n_cA^2)$ for the passes and the eigensubstitution
counts, plus one $O(A^3)$ back-transformation. The rescaling of $L_v$ and the
normalization of $\eta_v$ cancel in $P_{ab}$, because $Z_v$ is computed from the
same scaled vectors. $P$ may have one child, when the column is present in
only one of the two merged clades.

\section{An alignment-free initial time from trigram counts}
\label{app:kgram}

The E-step of Algorithm~\ref{alg:distance} needs an initial time $t_0$. A
single fixed value works well in practice, but $t_0$ can also be estimated
per pair, without alignment, from the number of shared trigrams,
$D_3=\sum_wc_X(w)\,c_Y(w)$, where $c_X(w)$ counts occurrences of the trigram
$w$ in $X$.

Given the alignment, residues in different columns are independent, so each
letter of a trigram matches with probability
$s(t)=\sum_a\pi_aM_{aa}(t)$ if the two positions are homologous and
$\rho=\sum_a\pi_a^2$ otherwise, under the LG chain used for distances. Taking $X$ as the ancestor, each
ancestral residue survives with probability $\alpha=e^{-\mu t}$. Fragments
are indivisible and insertions occur only between them, so adjacent residues
in one fragment stay adjacent, while across a fragment boundary a surviving
residue is followed directly by its surviving successor with probability
$(1-\beta)\alpha$. Following the fragtype of the current residue, the
probability that $j-1$ consecutive adjacencies are all preserved is
\[
p_j=\omega^\top\mathsf T^{\,j-1}\mathbf 1,\qquad
\mathsf T_{fg}=r_f\,\delta_{fg}+(1-r_f)(1-\beta)\alpha\,w_g,\qquad
\omega_f\propto\frac{w_f}{1-r_f}.
\]
Summing over the homology patterns of three consecutive diagonal position
pairs, and neglecting the rare pattern in which only the outer two are
homologous, gives, with $n=\alpha(L_X-2)$,
\[
\E[D_3]\approx\rho^3(L_X-2)(L_Y-2)
+n\Big(p_3(s^3-\rho^3)+2(p_2-p_3)(s^2\rho-\rho^3)
+(3-4p_2+p_3)(s\rho^2-\rho^3)\Big).
\]
Every term decreases in $t$, so $t_0$ is the root of $\E[D_3]=D_3$, found by
bisection in $O(L_X+L_Y+A^3)$ time per pair. Replacing $\rho$ by the pair's
own composition, $\sum_af_X(a)f_Y(a)$, corrects for composition bias. On
sequences simulated from the model, this estimate is accurate to within a
few percent up to $t\approx1.5$ and saturates beyond $t\approx2$; on real
protein families it is biased low. It is therefore suitable only as a
starting point for the EM iteration, not as a distance in its own right.

\begin{thebibliography}{9}
\bibitem[TKF-DP supplement(2026)]{HolmesMixFrag}
TKF-DP supplement, Appendix~A.3, ``The TKF92 model with fragment mixtures
(MixFrag)''. \url{https://tkfdp.net}.
\bibitem[Hobolth and Jensen(2005)]{HobolthJensen2005}
A.~Hobolth and J.~L. Jensen. Statistical inference in evolutionary models of
DNA sequences via the EM algorithm. \emph{Stat.\ Appl.\ Genet.\ Mol.\ Biol.}
4(1):18, 2005.
\bibitem[Holmes and Rubin(2002)]{HolmesRubin2002}
I.~Holmes and G.~M. Rubin. An expectation maximization algorithm for training
hidden substitution models. \emph{J.\ Mol.\ Biol.} 317(5):753--764, 2002.
\bibitem[Felsenstein(1981)]{Felsenstein81}
J.~Felsenstein. Evolutionary trees from {DNA} sequences: a maximum likelihood
approach. \emph{J.\ Mol.\ Evol.} 17:368--376, 1981.
\bibitem[Lartillot and Philippe(2004)]{Lartillot2004}
N.~Lartillot and H.~Philippe. A Bayesian mixture model for across-site
heterogeneities in the amino-acid replacement process. \emph{Mol.\ Biol.\
Evol.} 21:1095--1109, 2004.
\bibitem[Le et~al.(2008)]{LeGascuelLartillot2008}
S.~Q. Le, O.~Gascuel and N.~Lartillot. Empirical profile mixture models for
phylogenetic reconstruction. \emph{Bioinformatics} 24(20):2317--2323, 2008.
\bibitem[Le and Gascuel(2008)]{LeGascuel2008}
S.~Q. Le and O.~Gascuel. An improved general amino acid replacement matrix.
\emph{Mol.\ Biol.\ Evol.} 25(7):1307--1320, 2008.
\bibitem[Higgins and Sharp(1988)]{HigginsSharp1988}
D.~G. Higgins and P.~M. Sharp. CLUSTAL: a package for performing multiple
sequence alignment on a microcomputer. \emph{Gene} 73(1):237--244, 1988.
\bibitem[Thompson et~al.(1994)]{Thompson1994}
J.~D. Thompson, D.~G. Higgins and T.~J. Gibson. CLUSTAL W: improving the
sensitivity of progressive multiple sequence alignment through sequence
weighting, position-specific gap penalties and weight matrix choice.
\emph{Nucleic Acids Res.} 22(22):4673--4680, 1994.
\bibitem[Bradley et~al.(2009)]{Bradley2009}
R.~K. Bradley, A.~Roberts, M.~Smoot, S.~Juvekar, J.~Do, C.~Dewey, I.~Holmes
and L.~Pachter. Fast statistical alignment. \emph{PLoS Comput.\ Biol.}
5(5):e1000392, 2009.
\bibitem[Edgar(2004)]{Edgar2004}
R.~C. Edgar. MUSCLE: multiple sequence alignment with high accuracy and high
throughput. \emph{Nucleic Acids Res.} 32(5):1792--1797, 2004.
\bibitem[Do et~al.(2005)]{Do2005}
C.~B. Do, M.~S.~P. Mahabhashyam, M.~Brudno and S.~Batzoglou. ProbCons:
probabilistic consistency-based multiple sequence alignment. \emph{Genome
Res.} 15(2):330--340, 2005.
\bibitem[Hein et~al.(2000)]{Hein2000}
J.~Hein, C.~Wiuf, B.~Knudsen, M.~B. M{\o}ller and G.~Wibling. Statistical
alignment: computational properties, homology testing and goodness-of-fit.
\emph{J.\ Mol.\ Biol.} 302:265--279, 2000.
\bibitem[Thorne et~al.(1991)]{ThorneKishinoFelsenstein1991}
J.~L. Thorne, H.~Kishino and J.~Felsenstein. An evolutionary model for
maximum likelihood alignment of DNA sequences. \emph{J.\ Mol.\ Evol.}
33(2):114--124, 1991.
\bibitem[Thorne et~al.(1992)]{ThorneKishinoFelsenstein1992}
J.~L. Thorne, H.~Kishino and J.~Felsenstein. Inching toward reality: an
improved likelihood model of sequence evolution. \emph{J.\ Mol.\ Evol.}
34(1):3--16, 1992.
\bibitem[Fitch(1971)]{Fitch1971}
W.~M. Fitch. Toward defining the course of evolution: minimum change for a
specific tree topology. \emph{Syst.\ Zool.} 20(4):406--416, 1971.
\bibitem[Gascuel(1997)]{Gascuel1997}
O.~Gascuel. BIONJ: an improved version of the NJ algorithm based on a simple
model of sequence data. \emph{Mol.\ Biol.\ Evol.} 14(7):685--695, 1997.
\bibitem[Brejov{\'a} et~al.(2007)]{Brejova2007}
B.~Brejov{\'a}, D.~G. Brown and T.~Vina{\v r}. The most probable annotation problem in HMMs
and its application to bioinformatics. \emph{J.\ Comput.\ Syst.\ Sci.} 73(7):1060--1077, 2007.
\bibitem[Holmes and Durbin(1998)]{HolmesDurbin1998}
I.~Holmes and R.~Durbin. Dynamic programming alignment accuracy. \emph{J.\ Comput.\ Biol.}
5(3):493--504, 1998.
\bibitem[Lyngs{\o} and Pedersen(2002)]{LyngsoPedersen2002}
R.~B. Lyngs{\o} and C.~N.~S. Pedersen. The consensus string problem and the complexity of
comparing hidden Markov models. \emph{J.\ Comput.\ Syst.\ Sci.} 65(3):545--569, 2002.
\bibitem[Sankoff(1975)]{Sankoff1975}
D.~Sankoff. Minimal mutation trees of sequences. \emph{SIAM J.\ Appl.\ Math.}
28(1):35--42, 1975.
\end{thebibliography}

\end{document}
